\scanpage{251}
\mspage{126}

\uline{21 Lien avec les espaces de Teichmüller.}

D'abord un complément lié au diagramme \inserted{d'inclusions} (31) du n° 19, dans lequel \(A_{g,(\nu,\nu+1)}\) mérite la notation plus simple \(A_{g,\nu}^{\bullet}\) (\unsure{groupe} des autom. de \(U_{g,\nu}\) ponctué par \(a_{g,\nu}\)). Considérons \unsure{l'inclusion|la filtration} des sous-groupes de \(A_g\)
\[
 A_{g,\nu}^{\circ}\supset B_{g,\nu}\supset A_{g,\nu+1}^{\circ}
 \tag{72}
\]
\deleted{\unsure{\(A_{g,\nu}^{\bullet}\)}}

\unsure{donnant}, en divisant par \(A_{g,\nu+1}\) et dans le cas anabélien \(2g+\nu\geq 3\), la fibration

\begin{samepage}
(73) \diagram{\begin{tikzcd}[ampersand replacement=\&]
\widetilde U_{g,\nu} \arrow[r] \& U_{g,\nu} \\
A_{g,\nu}^{\circ}/A_{g,\nu+1}^{\circ} \arrow[u, phantom, "\unsure{\simeq}" description] \&
\end{tikzcd}}
\end{samepage}

On cherche le plus grand ss-groupe de \(A_g\) qui normalise \unsure{ces|les} triplet (72), et qui \deleted{\unsure{définit}} \inserted{opère} donc \deleted{\unsure{sur}} la fibration \(\widetilde U_{g,\nu}\to U_{g,\nu}\) (où \(\widetilde U_{g,\nu}\) est le rev. universel de \(U_{g,\nu}\) basé en \(a_{g,\nu}\)); comme les normalisateurs de \(A_{g,\nu}^{\circ}\) et \(A_{g,\nu+1}^{\circ}\) respectivement sont \(A_{g,\nu}\) et \(A_{g,\nu+1}\), le groupe en question doit être contenu dans leur intersection \(A_{g,\nu}^{\bullet}\), lequel normalise \deleted{\unsure{également}} \(B_{g,\nu}=A_{g,\nu}^{\bullet}\cap A_{g,\nu}^{\circ}\). \unsure{C'est} donc ce groupe qu'il y a lieu de faire opérer sur la fibration (73) -- le sous-groupe \unsure{formé} \deleted{\illegible{7} \(U_{g,\nu}\)} \inserted{des éléments de \(A_{g,\nu}\)} dont l'opération

\scanpage{252}

sur \(U_{g,\nu}\) est isotope à l'identité étant \unsure{justement} \(B_{g,\nu}\) -- donc c'est le groupe \unsure{\(\Gamma\)}
\[
 A_{g,\nu}^{\bullet}/B_{g,\nu}\simeq\Gamma_{g,\nu}
\]
qui, comme de juste, opère à isotopie près. On voudrait décrire le groupe de tous les autom. de la fibration topologique \(\widetilde U_{g,\nu}\) sur \(U_{g,\nu}\) qui sera a priori une extension du groupe
\[
 \begin{aligned}
 A_{g,\nu}&=\operatorname{Aut}(U_{g,\nu})\quad\text{par}\quad
 \Pi_{g,\nu}=\pi_1(\unsure{U_{g,\nu}|\widetilde U_{g,\nu}},a_{g,\nu})\\
 &=\operatorname{Aut}_{U_{g,\nu}}(\widetilde U_{g,\nu}).
 \end{aligned}
\]

\deleted{\illegible{1}} Cette extension \inserted{est \uline{\unsure{scindée can\textsuperscript{t}}}} \deleted{\unsure{splittée}} sur \(A_{g,\nu}^{\bullet}\), \inserted{stabilisateur au pt base \(a_{g,\nu}\)} (et on retrouve ainsi l'op. de \(A_{g,\nu}^{\bullet}\) pour la fibration (73). Or on a déjà une extension \(\Gamma'_{g,\nu+1}\) de \(\Gamma_{g,\nu}\) par \(\Pi_{g,\nu}\), d'où par image inverse par \(A_{g,\nu}\to\Gamma_{g,\nu}=A_{g,\nu}/A_{g,\nu}^{\circ}\), une extension (que je vais noter \(\widetilde A_{g,\unsure{\nu}}\)) de \(A_{g,\nu}\) par \(\Pi_{g,\nu}\) :
\[
 \widetilde A_{g,\nu}\simeq A_{g,\nu}\mathbin{\times}_{\Gamma_{g,\nu}}\Gamma'_{g,\nu+1}
 \tag{74}
\]
\(=\) groupe quotient du ss-groupe \deleted{\unsure{\(H\)}} \(\widetilde A_{g,\nu}^{\unsure{\natural|\sharp}}\) \unsure{de} \(A_{g,\nu}\times A_{g,\nu}^{\bullet}\) formé des couples \((u,v)\), \(u\unsure{\equiv}v\) \deleted{\unsure{\(A_{g,\nu}^{\circ}\)}} mod. \(A_{g,\nu}^{\circ}\), par le sous-groupe \deleted{\unsure{\(B_{g,\nu}^{\circ}\)}} \(1\times A_{g,\nu+1}^{\circ}\).

On a donc deux \inserted{(et = trois!)} structures d'extensions sur \(\widetilde A_{g,\nu}\)

\begin{samepage}
(75) \diagram{\begin{tikzcd}[ampersand replacement=\&,column sep=small]
1 \arrow[r] \& \Pi_{g,\nu} \arrow[r] \& \widetilde A_{g,\nu} \arrow[r,"p"] \& A_{g,\nu} \arrow[r] \& 1 \\
1 \arrow[r] \& A_{g,\nu}^{\circ} \arrow[r] \& \widetilde A_{g,\nu} \arrow[r] \& \Gamma'_{g,\nu+1} \arrow[r] \& 1
\end{tikzcd}}
\end{samepage}
\diagram{\begin{tikzcd}[ampersand replacement=\&,column sep=small]
1 \arrow[r] \& (A_{g,\nu}^{\circ}\unsure{\times|\rtimes}\Pi_{g,\nu}) \arrow[r] \& \widetilde A_{g,\nu} \arrow[r] \& \Gamma_{g,\nu} \arrow[r] \& 1
\end{tikzcd}}

\margin{\uline{NB.} Le deuxième \unsure{veut exprimer} que \(\widetilde A_{g,\nu}^{\circ}=A_{g,\nu}^{\circ}\) et \(\pi_0(\widetilde A_{g,\nu})\simeq\Gamma'_{g,\nu+1}\) (l'ex. can. de \(\Gamma_{g,\nu}\) par \(\Pi_{g,\nu}\)).}
\deleted{\unsure{Or}}

\scanpage{253}
\mspage{127}

Je voudrais faire opérer \(\widetilde A_{g,\nu}\) sur \(\widetilde U_{g,\nu}\), i.e. faire \deleted{\illegible{1}} opérer \(H=\widetilde A_{g,\nu}^{\unsure{\natural|\sharp}}\) \unsure{en opérat.} \deleted{de façon que \((u,v)\)} triviale de \(A_{g,\nu+1}^{\circ}\), et ceci en respectant les conditions suivantes

\margin{\unsure{donnée triviale} sur \(1\times A_{g,\nu+1}^{\circ}\). Cf généralisation ci-dessous en marge.}

a) \deleted{\illegible{1}} (compatibilité avec \(\widetilde A_{g,\nu}\to A_{g,\nu}\)) le couple \((u,v)\in H\) opère \unsure{sur} \(\widetilde U_{g,\nu}\) par un autom. compatible avec l'automorphisme \inserted{\(u\) de} \(U_{g,\nu}\) (on dira que c'est un \(u\)-autom.).

b) (compatibilité avec \(\Pi_{g,\nu}\to\widetilde A_{g,\nu}\)) Si \(u=1\), \inserted{[donc \(v\in A_{g,\nu}^{\circ}\), donc (comme \(v\in A_{g,\nu}^{\bullet}\)) \(v\in B_{g,\nu}\)]} alors l'opération de \((u,v)=(1,v)\) sur \(\widetilde U_{g,\nu}\) n'est autre que celle définie par \(v\bmod B_{g,\nu}^{\circ}=A_{g,\nu+1}^{\circ}\), qui est un élément de \(\Pi_{g,\nu}\simeq\operatorname{Aut}(\widetilde U_{g,\nu}/U_{g,\nu})\), ou encore, c'est celle définie par \uline{translation} \inserted{\uuline{à droite}} par \(\unsure{v^{-1}}\).

\margin{donnée de l'op. sur \(1\times B_{g,\nu}\).}

c) Compatibilité avec \deleted{\unsure{\(\widetilde A_{g,\nu}\to\Gamma'_{g,\nu+1}\)}} l'opération déjà \deleted{\illegible{1}} \inserted{obtenue plus haut} de \(A_{g,\nu}^{\bullet}\) sur \(\widetilde U_{g,\nu}\) : si \((u,v)\in H\) est tel que \(u=v\) (donc \(u=v\in A_{g,\nu}^{\bullet}\)) alors \(\unsure{(u,v)}=(u,u)\) opère \unsure{via} l'autom. int. défini par \(u\).

\margin{donnée de l'opération sur \(\delta(A_{g,\nu}^{\bullet})\) (\(\delta\) : flèche diagonale).}

d) \inserted{Compatibilité avec \(A_{g,\nu}^{\circ}\to\widetilde A_{g,\nu}\) :} L'opération de \(A_{g,\nu}^{\circ}\) sur \(\widetilde U_{g,\nu}\) est continue (ce qui, joint à a), détermine cette opération de manière unique \deleted{\illegible{2}} \inserted{et} l'existence \inserted{\unsure{a priori d'une telle op.}} (résulte de \(\pi_1(A_{g,\nu}^{\circ})=0\))).

\margin{donnée de l'opération sur \((A_{g,\nu}^{\circ})\times 1\).}

\unsure{tt} élément \inserted{\((u,v)\)} de \(H\) s'écrit de manière unique comme produit d'un élément \inserted{\((v,v)\)} de \(\delta(A_{g,\nu}^{\bullet})\), par un élément \inserted{\((v^{-1}u,1)\)} de \(A_{g,\nu}^{\circ}\times 1\) (le premier groupe normalisant le deuxième)

\margin{\uline{NB} L'opération de \(A_{g,\nu}^{\circ}\times 1\) se \unsure{décrit} le plus simplement par \uline{translations} de \unsure{\(A_{g,\nu}^{\circ}\)} sur l'espace \unsure{homogène} \(\widetilde U_{g,\nu}\) de \(A_{g,\nu}^{\circ}\).}

\scanpage{254}

et une opération de ce produit \(\frac12\) direct est donnée \inserted{bel et bien} par \deleted{\illegible{1}} la donnée d'opérations des deux groupes facteurs, satisfaisant une condition de compatibilité, qui est vérifiée ici par transport de structure. On a bien défini une opération de \(H=\widetilde A_{g,\nu}^{\unsure{\natural|\sharp}}\) sur la fibration (73), \deleted{\illegible{2}} satisfaisant les conditions c) et d) par construction même, il reste à vérifier a) et b).

Or pour a), il suffit de la vérifier séparément pour les \unsure{éléts} de \(H\) de \(\delta(A_{g,\nu}^{\bullet})\) et de \(A_{g,\nu}^{\circ}\times 1\), où c'est trivial par construction dans les deux cas. Reste à vérifier b), i.e. à expliciter l'opération d'un \unsure{élmt} \((1,v)\), \(v\in B_{g,\nu}\), qu'on écrit comme
\[
 (\qquad (1,v)=\deleted{(\unsure{v^{-1}},1)}(v,v)(v^{-1},1)
\]
donc
\[
 \rho\deleted{_{\unsure{\widetilde A}_{g,\nu}}}(1,v)=\rho(v,v)\rho(v^{-1},1).
\]

Or \(\rho(v,v)\) est induit par \(\operatorname{int}(v)\), \inserted{i.e. \(x\mapsto vxv^{-1}\)} et \(\rho(v^{-1},1)\) par translation à gauche : \(x\mapsto v^{-1}x\), donc le composé opère par \(x\mapsto xv^{-1}\). Ouf !

\scanpage{255}
\mspage{128}

On suppose maintenant que \(X_g\) est muni d'une structure \(C^\infty\), et on remplace dans les considérations \deleted{Soit} précédentes les groupes d'homéomorphismes par des groupes de difféomorphismes.

\margin{\uline{NB} Structure \(C^{\unsure{\nu|r}}\) avec \(\unsure{\nu|r}\geq 2\) suffirait}
\[
 \begin{aligned}
 E_g={}&\text{ens. des structures \deleted{complexes} \inserted{conformes}}\\
 &\text{sur }X_g\text{ compatibles avec sa}\\
 &\text{structure }C^\infty
 \end{aligned}
 \tag{76}
\]

\deleted{\unsure{Alors}} On voit de suite que \(E_g\) est un espace topologique \uline{\(\infty\)-connexe}, comme quotient de l'espace \(\infty\)-connexe \inserted{(un convexe)} des structures riemanniennes, par le \deleted{\unsure{gr.}} groupe \(\infty\)-connexe des applications \(C^\infty\) de \(X_g\) dans \(\mathbb R^{+*}\). Sur \(E_g\) le groupe \(A_g\) opère -- mais bien sûr \(E_g\) n'est plus un espace homogène. Notons tout de suite
\[
 \begin{aligned}
 E_g/A_g\simeq{}&\text{ens. des classes d'isomorphie}\\
 &\text{des \deleted{courbes algébriques} \inserted{surfaces} \deleted{\unsure{compactes}}}\\
 &\text{conformes compactes orientables}\\
 &\text{de genre }g
 \end{aligned}
 \tag{77}
\]

Et si on choisit une des deux orientations de \(X_g\), de sorte que \(E_g\) s'identifie à l'ens des structures \uline{complexes} sur \(X_g\) compatibles avec sa structure \(C^\infty\) \inserted{et son orientation}, alors \(E_g/A_g\) s'identifie à l'espace des classes d'isomorphie de courbes complexes \inserted{(non singulières) connexes} compactes de genre \(g\), \uline{modulo} passage à la complexe conjuguée. D'autre part

\scanpage{256}

\[
 \begin{aligned}
 E_g/A_g^+\simeq{}&\text{ens. des classes d'iso.}\\
 &\text{de courbes }\mathbb C\text{-alg \inserted{lisses} projectives}\\
 &\text{connexes de genre }g.
 \end{aligned}
 \tag{78}
\]

Plus généralement

\[
 \begin{aligned}
 E_g/A_{g,\nu}\simeq{}&\text{ens. des classes d'isomorphie}\\
 &\text{des \deleted{structures} \deleted{sur les} surfaces compactes}\\
 &\text{\inserted{conformes} munies \unsure{de multiponctures|de multiponctuations}}\\
 &\text{de type }(g,\nu)
 \end{aligned}
 \tag{79}
\]

\[
 \begin{aligned}
 E_g/A_{g,\nu}^{+}\simeq{}&\text{ens. des classes d'iso.}\\
 &\text{des courbes \deleted{\unsure{complexes}} \inserted{\(\mathbb C\) algébriques}}\\
 &\text{[lisses projectives connexes}\\
 &\text{de genre }g\text{] \uline{munies} d'une}\\
 &\text{multiponctuation, de type}\\
 &(g,\nu)
 \end{aligned}
 \tag{80}
\]
\margin{si \(X_g\) orientée.}

Ce sont là les ``espaces modulaires grossiers'' (``coarse moduli'' de Mumford), qu'on peut noter \(M_g^{\unsure{\natural|\sharp}}\) et \(M_{g,\nu}^{\unsure{\natural|\sharp}}\), et qui sont justement trop grossiers pour les usages géométriques les plus importants. \deleted{Mais} Beaucoup plus intéressants sont les quotients \(E_g/A_{g,\nu}^{\circ}\), \illegible{1}

\scanpage{257}
\mspage{129}

\[
 \widetilde M_{g,\nu}=E_g/A_{g,\nu}^{\circ}
 \qquad
 (\widetilde M_g=\widetilde M_{g,0}=E_g/A_g^{\circ})
 \tag{81}
\]
sur lesquels opère le groupe \deleted{\illegible{1}}
\[
 \Gamma_{g,\nu}=A_{g,\nu}/A_{g,\nu}^{\circ}
\]
par passage au quotient de l'op. de \(A_{g,\nu}\).

\deleted{\illegible{1}} L'espace \(\widetilde M_{g,\nu}\) \inserted{(avec l'opération de \(\Gamma_{g,\nu}\))} est appelé \uline{l'espace de Teichmüller} de type \(g,\nu\). Bien sûr on retrouve \(M_{g,\nu}^{\unsure{\natural|\sharp}}\) à partir de \(\widetilde M_{g,\nu}\) et \margin{\deleted{(82)}} de l'opération de \(\Gamma_{g,\nu}\) dessus, par
\[
 M_{g,\nu}^{\unsure{\natural|\sharp}}=\widetilde M_{g,\nu}/\Gamma_{g,\nu}
 \tag{\unsure{82}}
\]

\uline{Théorème} (Teichmüller) L'espace de Teichmüller \(\widetilde M_{g,\nu}\) est homéomorphe à \deleted{\(\mathbb R\)} \(\mathbb C^\mu\), où \(\mu=3g-3+\nu\) (\inserted{dans les} cas anabélien \(2g+\nu\geq 3\)), et \(\mu=3g-3+(\unsure{\nu}+\delta)\)
\[
 \delta=\dim_{\mathbb C}G
\]
dans le cas général, \(G\) étant \unsure{groupe} des autom. algébriques d'une \(U_{g,\nu}\) complexe (donc \(\delta=1\) dans les cas \unsure{``limites''} \((1,0)\) et \((0,2)\), et \inserted{plus \unsure{généralement}} \(\delta\) augmente de 1 chaque fois \unsure{pour} \(g\) fixé, et \((g,\nu-1)\) abélien, qu'on passe de \(\nu\) à \(\nu-1\)), donc

\margin{On peut dire aussi que \(\mu=3g-3+\nu'\), où \(\nu'=\nu\) dans le cas \unsure{anabélien}, et \unsure{où \(\nu'\)} dans le cas ``général'' est le plus petit entier \(\geq\nu\) tel que \((g,\nu')\) soit \unsure{``anabélien''}.}

\scanpage{258}

\[
 \left\{
 \begin{aligned}
 \mu&=1\quad\text{dans le cas }(1,0)\\
 \mu&=0\quad\text{dans les cas }(0,2),(0,1),(0,0),
 \end{aligned}
 \right.
 \tag{\unsure{\(82_1\)|\(82'\)}}
\]
i.e. \(\widetilde M_{g,\nu}\) est en ces cas réduit à un pt, i.e. l'action de \(A_{g,\nu}^{\circ}\) sur \(E_g=E_0\) est \deleted{déjà} transitive -- (et ce sont \inserted{avec le cas anabélien \((0,3)\)} les seuls trois cas où il en est ainsi, \deleted{\illegible{4}}).

\margin{i.e. sur surfaces \unsure{orientées} \unsure{sphériques à} \(\unsure{\nu\leq 3}\), toutes les structures complexes sont \uline{isotopes}.}

À partir de \(\widetilde M_{0,3}\), ou de \(\widetilde M_{1,1}\) ou de \(\widetilde M_{g,0}=\widetilde M_g\) avec \(g\geq 1\), quand \(\nu\) augmente les ``dimensions complexes'' de \(\widetilde M_{g,\nu}\) augmentent d'autant (Par contre \(\widetilde M_{1,0}\unsure{\longleftarrow}\widetilde M_{1,1}\) est un homéomorphisme \deleted{\unsure{parce que}\ \(A_{1,0}^{\circ}\unsure{\simeq}A_{1,1}^{\circ}\) \illegible{1}}).

On peut préciser le th. de Teichmüller ainsi

\uline{Corollaire} Dans le cas \uline{anabélien}, \(A_{g,\nu}^{\circ}\) opère librement sur \(E_g\), de sorte que \(E_g\) devient un torseur sur \(\widetilde M_{g,\nu}\), de groupe structural \(A_{g,\nu}^{\circ}\).

On en déduit que \(\widetilde M_{g,\nu}\) joue le rôle d'un \uline{espace classifiant} pour \(A_{g,\nu}^{\circ}\), et que
\[
 \pi_{\unsure{i|\nu}}(\widetilde M_{g,\nu})
 \xrightarrow{\sim} \pi_{\unsure{\nu-1|i-1}}(A_{g,\nu}^{\circ})
 \tag{83}
\]
et le fait que \(\widetilde M_{g,\nu}\) soit \(\infty\)-connexe (contenu dans le th. de Teichmüller) équivaut à

\scanpage{259}
\mspage{130}

celui que \(A_{g,\nu}^{\circ}\) le soit -- ce qui est le ``théorème bien connu'' rappelé au n° 19.

Je n'insiste pas ici sur l'interprétation des points de \(\widetilde M_{g,\nu}\) comme des classes d'isom. de courbes complexes, munies \deleted{de} \inserted{d'une} ``rigidification de Teichmüller'' convenable, \unsure{et le pt de vue espaces modulaires rigidifiés}, qui permet de vérifier a priori que \(\widetilde M_{g,\nu}\) est muni d'une structure de variété complexe non singulière -- mais déjà le fait que \(\widetilde M_{g,\nu}\) soit \uline{simplement connexe} (ce qui peut \unsure{exprimer} \unsure{en} interprétant \(\widetilde M_{g,\nu}\) comme rev. universel du \uline{topos modulaire} \(M_{g,\nu}\)) est bien profond, qui ne semble pas pouvoir rentrer dans le cadre de la topologie (ou de la topologie différentielle \((C^\infty)\) sans plus...

\margin{\uline{NB} Pour prouver que les \(\unsure{\widetilde M_{g,\nu}}\) sont \(\infty\)-connexes, \unsure{on est réduit facilement} au cas de \(M_{g,0}\) \unsure{(si \(g\geq 2\))}, ou de \(M_{1,1}\) (cas elliptique ponctué), en utilisant \unsure{les fibrations} \(A_{g,\nu}^{\circ}/A_{g,\nu+1}^{\circ}\simeq\widetilde U_{g,\nu}\) (qui est \(\infty\)-connexe...). Le cas \(g=1\) est \unsure{d'ailleurs facile et} bien compris...}

\scanpage{260}

Notons que les inclusions de sous-groupes
\[
 A_{g,\nu+1}^{\circ}\subset A_{g,\nu}^{\circ}\cdots\subset A_{g,0}^{\circ}=A_g^{\circ}
\]
définissent une tour d'applications continues
\[
 \cdots\longrightarrow\widetilde M_{g,\nu+1}\longrightarrow\widetilde M_{g,\nu}
 \longrightarrow\cdots\longrightarrow\widetilde M_{g,0}=\widetilde M_g
 \tag{84}
\]

où
\[
 \widetilde M_{g,\nu+1}\longrightarrow\widetilde M_{g,\nu}
 \tag{85}
\]
est \inserted{\uline{pour \((g,\nu)\) anabéliens}} une fibration en fibre \(A_{g,\nu}^{\circ}/A_{g,\nu+1}^{\circ}\simeq\widetilde U_{g,\nu}\).

[il est donc à fibre contractile, et la \(\infty\)-connexité de \(\widetilde M_{g,\nu+1}\) équivaut donc à celle de \(\widetilde M_{g,\nu}\) -- ce qui nous ramène au cas de \(\widetilde M_g\) si \(g\geq 2\), de \(M_{\text{\deleted{\(g\)}}1,1}\) et de \(M_{0,3}\) si \(g\unsure{=|\neq|\geq}1\) ou 0].

Dans le cas \((g,\nu)\) abélien, on trouve
\[
 \widetilde M_{g,\nu+1}\simeq\widetilde M_{g,\nu}\qquad
 (\text{si }(g,\nu)\text{ abélien})
 \tag{85}
\]
ce qui \deleted{\unsure{précisément résulte|provient, résulte} du fait déjà} \inserted{signifie} \deleted{signalé} que dans ces cas, si on a deux structures complexes \inserted{\(\alpha,\beta\)} sur \(X\) qui sont congrues par \inserted{\(u\in\)} \(A_{g,\nu}^{\circ}\), \inserted{(i.e. \(\beta=u.\alpha\))} elles sont = congrues par \(A_{g,\nu+1}^{\circ}\).

En effet, si \(G\) est \unsure{la} \inserted{composante \unsure{unité|neutre} du} groupe des autom. complexes de \(U_{g,\nu}\) pour \(\beta\), on peut écrire \(u=gv\) avec \(g\in G\), \(v\in A_{g,\nu+1}^{\circ}\) donc \(\beta=gv.\alpha\) donc (comme \(g^{-1}\beta=\beta\)) \(\beta=v\alpha\), cqfd.
